% =====================================================================
%  Capítulo 2 — Acceso al panel
% =====================================================================
\chapter{Acceso al panel}\label{cap:acceso}

\entradilla{Cómo entrar al panel, recuperar la contraseña y moverse por la interfaz. Lo que
aprenda en este capítulo le sirve para todos los módulos.}

\section{Iniciar sesión}\label{sec:iniciar-sesion}

\ficha{F084}{Pantalla de acceso de Trainova}{/login}{Entrar al panel con su correo y contraseña.}

\begin{pasos}
  \item Abra el navegador y escriba la dirección de Trainova seguida de \ruta{/login}, es decir,
        \ruta{\dominio\allowbreak/login}.
        También puede entrar desde el sitio público: al final de la página de inicio está el
        enlace \textbf{Acceso administradores}.
  \item Verá el logo de Training Corporation y el título \textbf{Panel de cotizaciones}.
  \item En \campo{Correo electrónico}, escriba el correo de su cuenta.
  \item En \campo{Contraseña}, escriba su contraseña.
  \item Si está en su equipo personal, puede marcar \campo{Mantener la sesión iniciada}. Así el
        navegador conserva la sesión abierta aunque lo cierre, hasta que usted cierre sesión.
  \item Haga clic en \boton{Iniciar sesión}.
\end{pasos}

\captura{02-01-login}{/login}{Pantalla de inicio de sesión.}
  {la pantalla de inicio de sesión completa, con el logo, el título «Panel de cotizaciones»
   y los campos vacíos.}

\begin{resultado}
El sistema lo lleva al \textbf{Panel de inicio} (\ruta{/dashboard}), donde lo saluda por su
nombre. Continúe en la \seccion{sec:panel-inicio}.
\end{resultado}

\begin{importante}
No marque \campo{Mantener la sesión iniciada} en computadores compartidos o públicos:
cualquier persona que use ese navegador después podría entrar al panel con su cuenta.
\end{importante}

\subsection{Si no puede entrar}\label{sec:errores-acceso}

Cuando algo falla, el formulario muestra un mensaje en rojo debajo del campo afectado.

\begin{tabla}{colspec={Q[3.6cm,l] X[l] Q[4.4cm,l]}}
  Situación & Qué verá & Qué hacer \\
  Dejó un campo vacío & \msg{El correo electrónico es obligatorio.} o
    \msg{La contraseña es obligatoria.} & Complete el campo y vuelva a intentar. \\
  El correo no tiene un formato válido & El navegador le pide corregirlo, o el sistema muestra
    \msg{Ingresa un correo electrónico válido.} & Revise que el correo tenga la forma
    \texttt{nombre@empresa.com}. \\
  El correo o la contraseña no coinciden & Un mensaje bajo el correo. Actualmente aparece en
    inglés: \msg{These credentials do not match our records.} & Revise mayúsculas y la tecla
    \tecla{Bloq Mayús}. Si no la recuerda, recupérela (\seccion{sec:recuperar}). \\
  Varios intentos fallidos seguidos & Tras 5 intentos fallidos el acceso se bloquea por unos
    segundos. El mensaje aparece en inglés: \msg{Too many login attempts. Please try again in
    \var{N} seconds.} & Espere el tiempo indicado e intente de nuevo. \\
\end{tabla}

\captura{02-02-login-error}{/login}{Mensaje de error al iniciar sesión.}
  {el formulario de inicio de sesión después de escribir una contraseña incorrecta, con el
   mensaje en rojo visible debajo del correo.}

\section{Recuperar la contraseña}\label{sec:recuperar}

Si olvidó su contraseña, puede crear una nueva desde la pantalla de inicio de sesión.

\begin{pasos}
  \item En la pantalla de inicio de sesión, haga clic en \textbf{¿Olvidó su contraseña?}
  \item Verá el texto \msg{¿Olvidó su contraseña? Escriba su correo electrónico y le enviaremos
        un enlace para crear una nueva.} Escriba en el campo el correo de su cuenta.
  \item Haga clic en \boton{Enviar enlace para restablecer contraseña}. Verá un mensaje de
        confirmación en verde (actualmente en inglés).
  \item Revise su bandeja de entrada y, si no lo encuentra, la carpeta de correo no deseado.
        El correo llega en inglés, con el asunto \msg{Reset Password Notification}.
  \item Abra el correo y haga clic en el botón \textbf{Reset Password}.
  \item En la pantalla \textbf{Restablecer contraseña}, verifique su correo en
        \campo{Correo electrónico}, escriba la nueva contraseña en \campo{Contraseña} y repítala
        en \campo{Confirmar contraseña}. Debe tener al menos 8 caracteres.
  \item Haga clic en \boton{Reset Password}.
\end{pasos}

\captura{02-03-recuperar-contrasena}{/forgot-password}{Solicitud del enlace para restablecer la contraseña.}
  {la pantalla «¿Olvidó su contraseña?» con el campo de correo y el botón «Enviar enlace para
   restablecer contraseña».}

\begin{resultado}
Vuelve a la pantalla de inicio de sesión con un mensaje de confirmación. Inicie sesión con la
nueva contraseña.
\end{resultado}

\begin{nota}
El enlace del correo vence a los 60 minutos y solo puede pedirse uno por minuto. Si vence,
solicite uno nuevo. Otra opción es pedirle a otro administrador que le asigne una contraseña
nueva desde \menu{Sistema}{Usuarios} (ver \seccion{sec:usuario-editar}).
\end{nota}

\section{El Panel de inicio}\label{sec:panel-inicio}

\ficha{F015}{Primera pantalla después de iniciar sesión}{/dashboard}{Ir con un clic a cualquier
módulo del panel.}

Después de iniciar sesión verá el saludo \textbf{Bienvenido/a, \var{su primer nombre}} y el texto
\msg{Gestione las solicitudes de cotización, el catálogo que ve el cliente y la configuración
del sistema.} Debajo aparecen tarjetas de acceso directo, agrupadas igual que la barra lateral.
Haga clic en una tarjeta para abrir el módulo.

\captura[8.5cm]{02-04-panel-inicio}{/dashboard}{Panel de inicio del administrador.}
  {la pantalla completa del Panel justo después de iniciar sesión: la barra lateral azul a la
   izquierda, el saludo «Bienvenido/a, …» y todas las tarjetas de los grupos Operación, Catálogo y
   Sistema.}

\begin{tabla}{colspec={Q[2.2cm,l] Q[3.8cm,l] X[l] Q[1.9cm,c]}}
  Grupo & Tarjeta & Descripción en pantalla & Sección \\
  Operación & \icono{\icCotizacion}\ Cotizaciones & Solicitudes recibidas y su estado & \ref*{sec:cotizaciones} \\
  Operación & \icono{\icContactos}\ Clientes & Quién ha cotizado y su historial & \ref*{sec:clientes} \\
  Catálogo & \icono{\icMaletin}\ Unidades de negocio & Unidades y su código de cotización & \ref*{sec:unidades} \\
  Catálogo & \icono{\icRama}\ Líneas de gestión & Agrupan los servicios ofrecidos & \ref*{sec:lineas} \\
  Catálogo & \icono{\icLlave}\ Servicios & Servicios por unidad y línea & \ref*{sec:servicios} \\
  Catálogo & \icono{\icArbol}\ Árbol de decisión & Preguntas del cotizador & \ref*{sec:arbol} \\
  Sistema & \icono{\icUsuarios}\ Usuarios & Cuentas con acceso al panel & \ref*{sec:usuarios} \\
  Sistema & \icono{\icCorreo}\ Correos de bienvenida & Invitar clientes a la plataforma & \ref*{sec:bienvenida} \\
  Sistema & \icono{\icAjustes}\ Configuración & Notificaciones y datos de contacto & \ref*{sec:configuracion} \\
\end{tabla}

\begin{consejo}
Para volver al Panel de inicio desde cualquier pantalla, haga clic en el logo o en el texto
\textbf{Panel de cotizaciones}, en la parte superior de la barra lateral.
\end{consejo}

\section{Partes de la pantalla}\label{sec:partes}

Todas las pantallas del panel comparten la misma estructura. La \figref{fig:esquema-pantalla}
la muestra en forma de esquema; las letras se explican a continuación.

\begin{figure}[H]
\centering
\begin{tikzpicture}[x=1cm, y=1cm,
    menu/.style={anchor=west, inner sep=0pt, text=white, font=\montserrat\fontseries{sb}\selectfont\fontsize{6}{7}\selectfont},
    grupomenu/.style={anchor=west, inner sep=0pt, text=white!80!azul, font=\montserrat\bfseries\fontsize{5.5}{6}\selectfont},
    marca/.style={circle, fill=noche, draw=white, line width=1pt, text=white, inner sep=0pt,
      minimum size=0.52cm, font=\montserrat\bfseries\scriptsize},
    textofalso/.style={fill=linea, rounded corners=1pt},
  ]
  % Marco
  \draw[linea, line width=1pt, fill=nube, rounded corners=4pt] (0,0) rectangle (15.6,8.8);
  % Barra lateral
  \fill[azul, rounded corners=4pt] (0,0) rectangle (3.6,8.8);
  \fill[azul] (3.0,0) rectangle (3.6,8.8);
  \node[anchor=north west, inner sep=0pt] at (0.35,8.45) {\includegraphics[width=1.5cm]{imagenes/logo-blanco.png}};
  \node[menu, font=\montserrat\bfseries\fontsize{6}{7}\selectfont] at (0.35,7.45) {Panel de cotizaciones};
  \node[grupomenu] at (0.35,6.95) {Operación};
  \fill[white, rounded corners=1.5pt] (0.22,6.42) rectangle (3.38,6.74);
  \node[menu, text=azul] at (0.38,6.58) {\ico{F70E}\hspace{0.4em}Cotizaciones};
  \node[menu] at (0.38,6.18) {\ico{F2B9}\hspace{0.4em}Clientes};
  \node[grupomenu] at (0.35,5.7) {Catálogo};
  \node[menu] at (0.38,5.32) {\ico{F0B1}\hspace{0.4em}Unidades de negocio};
  \node[menu] at (0.38,4.98) {\ico{F126}\hspace{0.4em}Líneas de gestión};
  \node[menu] at (0.38,4.64) {\ico{F0AD}\hspace{0.4em}Servicios};
  \node[menu] at (0.38,4.30) {\ico{F0E8}\hspace{0.4em}Árbol de decisión};
  \node[grupomenu] at (0.35,3.82) {Sistema};
  \node[menu] at (0.38,3.44) {\ico{F0C0}\hspace{0.4em}Usuarios};
  \node[menu] at (0.38,3.10) {\ico{F0E0}\hspace{0.4em}Correos de bienvenida};
  \node[menu] at (0.38,2.76) {\ico{F013}\hspace{0.4em}Configuración};
  \draw[white!60!azul, line width=0.5pt] (0.22,1.35) -- (3.38,1.35);
  \node[menu] at (0.38,1.05) {\ico{F08E}\hspace{0.4em}Ver sitio público};
  \node[menu] at (0.38,0.62) {\ico{F2F5}\hspace{0.4em}Cerrar sesión};
  \node[menu, font=\fontsize{5}{6}\selectfont, text=white!80!azul] at (0.75,0.36) {usuario@empresa.com};
  % Aviso
  \draw[azul, fill=white, line width=0.8pt, rounded corners=2pt] (4.2,7.75) rectangle (15.0,8.3);
  \node[anchor=west, inner sep=0pt, font=\fontsize{6.5}{7}\selectfont, text=tinta] at (4.4,8.02)
    {{\color{azul}\ico{F058}}\hspace{0.5em}Servicio creado correctamente.};
  \node[anchor=east, inner sep=0pt, font=\fontsize{6.5}{7}\selectfont, text=pizarra] at (14.8,8.02) {\ico{F00D}};
  % Encabezado de la pantalla
  \node[anchor=west, inner sep=0pt, font=\montserrat\fontseries{eb}\selectfont\fontsize{13}{14}\selectfont, text=azul] at (4.2,7.2) {Título de la pantalla};
  \fill[textofalso] (4.2,6.62) rectangle (10.4,6.78);
  \node[fill=azul, text=white, font=\montserrat\bfseries\fontsize{6.5}{7}\selectfont, rounded corners=2pt,
    inner xsep=6pt, inner ysep=4pt, anchor=east] at (15.0,7.05) {\ico{2B}\hspace{0.4em}Nuevo};
  % Búsqueda
  \draw[grismarca, fill=white, line width=0.7pt, rounded corners=2pt] (4.2,5.75) rectangle (10.6,6.25);
  \node[anchor=west, inner sep=0pt, font=\fontsize{6.5}{7}\selectfont, text=pizarra] at (4.4,6.0) {\ico{F002}\hspace{0.5em}Buscar…};
  \draw[azul, line width=0.8pt, rounded corners=2pt] (10.75,5.75) rectangle (11.95,6.25);
  \node[font=\montserrat\bfseries\fontsize{6.5}{7}\selectfont, text=azul] at (11.35,6.0) {Buscar};
  % Tabla
  \draw[linea, fill=white, line width=0.8pt, rounded corners=3pt] (4.2,1.55) rectangle (15.0,5.4);
  \fill[nube] (4.25,4.85) rectangle (14.95,5.35);
  \foreach \y in {4.85, 4.15, 3.45, 2.75, 2.05} { \draw[linea, line width=0.5pt] (4.2,\y) -- (15.0,\y); }
  \foreach \y in {4.5, 3.8, 3.1, 2.4, 1.8} {
    \fill[textofalso] (4.5,\y-0.06) rectangle (5.0,\y+0.06);
    \fill[textofalso] (5.6,\y-0.06) rectangle (8.6,\y+0.06);
    \fill[textofalso] (9.2,\y-0.06) rectangle (11.4,\y+0.06);
  }
  \foreach \y in {4.5, 3.8, 3.1, 2.4, 1.8} {
    \node[font=\fontsize{6.5}{7}\selectfont, text=pizarra] at (13.2,\y) {\ico{F06E}};
    \node[font=\fontsize{6.5}{7}\selectfont, text=azul] at (13.75,\y) {\ico{F303}};
    \node[font=\fontsize{6.5}{7}\selectfont, text=error] at (14.3,\y) {\ico{F2ED}};
  }
  % Paginación y pie
  \node[anchor=west, inner sep=0pt, font=\fontsize{6}{7}\selectfont, text=pizarra] at (4.2,1.2) {Mostrando 1–15 de 40};
  \node[anchor=east, inner sep=0pt, font=\fontsize{6}{7}\selectfont, text=pizarra] at (15.0,1.2)
    {\ico{F053}\hspace{0.3em}Anterior\hspace{1.2em}1 / 3\hspace{1.2em}Siguiente\hspace{0.3em}\ico{F054}};
  \draw[linea, line width=0.6pt] (3.6,0.62) -- (15.6,0.62);
  \node[font=\fontsize{5.5}{6}\selectfont, text=pizarra] at (9.6,0.33) {© 2026 Training Corporation. Todos los derechos reservados.};
  % Marcas
  \node[marca] at (3.6,5.0) {A};
  \node[marca] at (3.6,1.05) {B};
  \node[marca] at (15.0,8.3) {C};
  \node[marca] at (15.6,6.95) {D};
  \node[marca] at (12.25,6.0) {E};
  \node[marca] at (15.0,3.45) {F};
  \node[marca] at (9.6,1.2) {G};
\end{tikzpicture}
\caption{Esquema de una pantalla del panel.}\label{fig:esquema-pantalla}
\end{figure}

\begin{zonas}
  \item \textbf{Barra lateral.} Contiene el logo, el texto \textbf{Panel de cotizaciones} y los
        módulos agrupados en \textbf{Operación}, \textbf{Catálogo} y \textbf{Sistema}. El módulo
        donde se encuentra aparece resaltado en blanco.
  \item \textbf{Pie de la barra lateral.} \textbf{Ver sitio público} abre Trainova en otra
        pestaña y \textbf{Cerrar sesión} termina su sesión. Debajo aparece el correo de la cuenta
        con la que entró.
  \item \textbf{Aviso del sistema.} Confirma el resultado de la última acción
        (\seccion{sec:avisos}).
  \item \textbf{Encabezado.} Título azul de la pantalla, una descripción en gris y, a la
        derecha, los botones de acción, como \boton{\icMas\ Nuevo servicio}.
  \item \textbf{Búsqueda y filtros} (\seccion{sec:busqueda}).
  \item \textbf{Tabla} con los registros y los íconos de acción de cada fila
        (\seccion{sec:acciones-fila}).
  \item \textbf{Paginación} y pie de página con el aviso de derechos de Training Corporation.
\end{zonas}

\capturaconnotas[11.5cm]{02-05-barra-lateral}{/dashboard}{Barra lateral del panel.}
  {un recorte vertical de la barra lateral azul completa, desde el logo hasta «Cerrar sesión»,
   con algún módulo seleccionado.}
  {\textbf{Cómo moverse por el panel}

   \begin{itemize}
     \item Haga clic en el nombre de un módulo para abrirlo.
     \item El módulo activo se resalta con fondo blanco y letra azul.
     \item Los grupos solo organizan los módulos: no se pueden abrir ni cerrar.
     \item Si la pantalla es baja, la barra lateral se desplaza sola con la rueda del ratón.
     \item \textbf{Ver sitio público} se abre en otra pestaña; el panel sigue abierto en la
           pestaña actual.
   \end{itemize}}

\section{Elementos comunes}\label{sec:comunes}

Las herramientas de esta sección se repiten en casi todos los módulos. Una vez las conozca,
podrá usarlas igual en cualquier pantalla.

\subsection{Avisos del sistema}\label{sec:avisos}

Después de guardar, eliminar o enviar algo, aparece un aviso en la parte superior del contenido:

\begin{itemize}
  \item \textbf{Éxito:} borde azul con un ícono \icono{\icOk}. Por ejemplo,
        \msg{Línea de gestión creada correctamente.}
  \item \textbf{Error:} recuadro rojo con un ícono de alerta. Por ejemplo,
        \msg{No puedes eliminar tu propio usuario.}
\end{itemize}

El aviso desaparece solo a los 8 segundos. Para cerrarlo antes, haga clic en la \icono[pizarra]{\icCerrar}
de la derecha.

\captura[3.2cm]{02-06-aviso}{/dashboard/lineas-gestion}{Aviso de confirmación después de guardar.}
  {un recorte de la parte superior del contenido justo después de crear o editar un registro,
   con el aviso azul de éxito visible.}

\subsection{Búsqueda}\label{sec:busqueda}

\begin{pasos}
  \item Escriba en la caja de búsqueda (la que tiene el ícono \icono[pizarra]{\icBuscar}). El texto gris
        de la caja indica qué datos puede buscar en esa pantalla.
  \item Presione \tecla{Enter} o haga clic en \botonsec{Buscar}.
  \item Para quitar la búsqueda, haga clic en la \icono[pizarra]{\icCerrar} \textbf{Limpiar búsqueda}
        que aparece junto a la caja.
\end{pasos}

\begin{nota}
La búsqueda encuentra textos parciales y no distingue mayúsculas de minúsculas. Si escribe
solo números, también busca por el número de identificación (ID) del registro. La búsqueda no
se ejecuta mientras escribe: debe presionar \tecla{Enter} o \botonsec{Buscar}.
\end{nota}

\subsection{Filtros}\label{sec:filtros}

Algunas pantallas tienen listas desplegables junto a la búsqueda, por ejemplo
\campo{Unidad de negocio} o \campo{Empresa}. El filtro se aplica en cuanto elige una opción.
Elija \textbf{Todas} para quitarlo.

\subsection{Tablas y paginación}\label{sec:paginacion}

Las tablas muestran 15 registros por página. Debajo de la tabla verá el texto
\msg{Mostrando \var{desde}–\var{hasta} de \var{total}} y, si hay más de una página, los botones
\botonsec{\icAnterior\ Anterior} y \botonsec{Siguiente\ \icSiguiente} con la página actual
(por ejemplo, \textbf{1 / 3}). Al cambiar la búsqueda o un filtro, la tabla vuelve a la
primera página.

\begin{consejo}
La dirección de la página guarda la búsqueda, los filtros y la página. Puede guardarla como
favorito en el navegador o enviársela a otro administrador para que vea exactamente la misma
lista.
\end{consejo}

En pantallas angostas, las tablas se desplazan hacia los lados. Una sombra en el borde indica
que hay más columnas por ver.

\subsection{Acciones de cada fila}\label{sec:acciones-fila}

La columna \campo{Acciones} muestra íconos. Si pasa el ratón sobre uno, aparece su nombre.

\begin{tabla}{colspec={Q[1.6cm,c,m] Q[3cm,l,m] X[l,m]}}
  Ícono & Nombre & Qué hace \\
  {\color{pizarra}\fontsize{12}{12}\selectfont\icVer} & Ver detalle & Abre la ficha completa del registro. \\
  {\color{azul}\fontsize{12}{12}\selectfont\icEditar} & Editar & Abre el formulario para modificar el registro. \\
  {\color{error}\fontsize{12}{12}\selectfont\icEliminar} & Eliminar & Borra el registro, previa confirmación. \\
  {\color{azul}\fontsize{12}{12}\selectfont\icExterno} & Abrir propuesta & Solo en Cotizaciones: abre la propuesta enviada en una pestaña nueva. \\
\end{tabla}

\captura{02-07-lista-tipica}{/dashboard/unidad-negocio}{Una lista típica: búsqueda, tabla, acciones y paginación.}
  {la pantalla de Unidades de negocio completa, con el encabezado, el botón «Nueva unidad»,
   la caja de búsqueda, la tabla y los íconos de acción de cada fila.}

\subsection{Confirmar una eliminación}\label{sec:confirmar}

Antes de borrar algo, el navegador muestra un cuadro de confirmación con un mensaje propio de
cada módulo, por ejemplo \msg{¿Eliminar el servicio «Trabajo en Alturas»?}

\begin{itemize}
  \item Haga clic en \textbf{Aceptar} para eliminar.
  \item Haga clic en \textbf{Cancelar} para conservar el registro.
\end{itemize}
Los botones los dibuja el navegador, no Trainova: en un navegador en español dicen
\textbf{Aceptar} y \textbf{Cancelar}; en uno en inglés, como el de la captura, dicen
\textbf{OK} y \textbf{Cancel}. El cuadro también muestra la dirección \ruta{\dominio}.

\captura[4.5cm]{02-08-confirmar-eliminar}{/dashboard/servicios}{Cuadro de confirmación antes de eliminar.}
  {el cuadro de confirmación del navegador que aparece al hacer clic en el ícono rojo de
   eliminar de cualquier fila. Haga clic en «Cancelar» después de capturarlo.}

\begin{importante}
Las eliminaciones no se pueden deshacer. Si borró algo por error, tendrá que crearlo de nuevo.
\end{importante}

\subsection{Formularios}\label{sec:formularios}

Los formularios para crear y editar registros funcionan igual en todo el panel:

\begin{itemize}
  \item Arriba a la izquierda hay un botón de regreso, \botonsec{\icAtras\ Volver a \var{módulo}}
        (por ejemplo, \textbf{Volver a cotizaciones}), que lleva a la lista.
  \item Los campos obligatorios no llevan marca. Los opcionales muestran
        \textbf{(opcional)} junto al nombre.
  \item El texto gris bajo el nombre de un campo explica qué escribir.
  \item Al guardar, el sistema revisa los datos. Si algo falta o está mal, el mensaje aparece en
        rojo debajo del campo y no se guarda nada hasta corregirlo.
  \item \botonsec{Cancelar} vuelve a la lista sin guardar.
  \item El botón principal, por ejemplo \boton{Crear servicio} o \boton{Guardar cambios},
        guarda los datos. Mientras procesa, muestra \textbf{Guardando…}.
\end{itemize}

\subsection{Fechas y valores vacíos}

Las fechas se muestran en formato día/mes/año (\texttt{27/10/2026}). Cuando un dato no fue
diligenciado, el panel muestra una raya (\textbf{—}).

\section{Usar el panel desde el celular}\label{sec:movil}

En pantallas angostas, como celulares o tabletas en vertical, la barra lateral se oculta para
dejar espacio al contenido.

\begin{pasos}
  \item Toque el botón \textbf{\icMenu} (\textbf{Abrir menú}), arriba a la izquierda, junto al
        texto \textbf{Panel de cotizaciones}.
  \item Se despliega la barra lateral. Toque el módulo que necesita.
  \item Para cerrarla sin elegir nada, toque la \textbf{\icCerrar} o cualquier punto fuera de la barra.
\end{pasos}

\begin{nota}
En pantallas angostas, algunas tablas ocultan columnas para que quepan. En
\menu{Operación}{Cotizaciones}, la unidad de negocio y la fecha pasan a mostrarse debajo del
código; en \menu{Operación}{Clientes} se ocultan las columnas \campo{Teléfono} y \campo{Última}.
Abra la ficha del registro para ver todos los datos.
\end{nota}

\capturamovil{02-09-menu-movil}{Menú del panel abierto en un celular.}
  {el panel en vista de celular con el menú lateral abierto. En Chrome puede simularlo con
   F12 y el ícono de dispositivo móvil (Ctrl + Shift + M).}
